% =====================================================================
%  DEEP LEARNING -- HOMEWORK 3-4 REPORT
%  "CNNs & Computer Vision with PyTorch" -- Weeks 3-4.
% =====================================================================
\documentclass[10pt,twocolumn,letterpaper]{article}

% ---- fonts & encoding (Times, IEEE-like) ----------------------------
\usepackage[T1]{fontenc}
\usepackage{mathptmx}

% ---- IEEE-like page geometry ----------------------------------------
\usepackage[letterpaper,top=0.75in,bottom=1in,left=0.625in,right=0.625in,
            columnsep=0.25in]{geometry}

% ---- math, graphics, tables, misc -----------------------------------
\usepackage{amsmath,amssymb}
\usepackage{graphicx}
\graphicspath{{figures/}{./}}
\usepackage{booktabs}
\usepackage{caption}
\usepackage{enumitem}
\usepackage{xcolor}
\usepackage{listings}
\usepackage{titlesec}
\usepackage{cite}
\usepackage[hidelinks]{hyperref}

% ---- IEEE-style section headings ------------------------------------
\renewcommand{\thesection}{\Roman{section}}
\renewcommand{\thesubsection}{\Alph{subsection}}
\titleformat{\section}[block]{\centering\normalsize\scshape}{\thesection.}{0.5em}{}
\titleformat{\subsection}[block]{\normalfont\itshape}{\thesubsection.}{0.5em}{}
\titlespacing*{\section}{0pt}{9pt}{4pt}
\titlespacing*{\subsection}{0pt}{6pt}{3pt}

% ---- IEEE-style captions -------------------------------------------
\renewcommand{\thetable}{\Roman{table}}
\captionsetup[figure]{name={Fig.},labelsep=period,font=small}
\captionsetup[table]{labelsep=newline,textfont=sc,font=small,justification=centering}

% ---- code listing style ---------------------------------------------
\definecolor{codekw}{RGB}{0,0,180}
\definecolor{codecm}{RGB}{110,110,110}
\lstset{
  language=Python, basicstyle=\ttfamily\footnotesize, keywordstyle=\color{codekw}\bfseries,
  commentstyle=\color{codecm}\itshape, showstringspaces=false, breaklines=true,
  columns=fullflexible, frame=single, framesep=4pt, aboveskip=6pt, belowskip=6pt,
}

\begin{document}

% =====================================================================
%  TITLE / AUTHOR BLOCK
% =====================================================================
\twocolumn[{%
\begin{center}
  {\LARGE\bfseries CNNs and Computer Vision with PyTorch:\\[2pt]
   Architecture, Transfer Learning, and Data Augmentation\par}
  \vspace{4pt}
  {\normalsize Deep Learning Cohort~I 2026 \textperiodcentered\ Homework~3--4 Report (Weeks~3--4)\par}
  \vspace{10pt}
  {\large Bayram Bayramov\par}
  \vspace{2pt}
  {\normalsize Student ID: 121314 \quad Group: M003\par}
  {\normalsize bayramovb578@gmail.com\par}
  \vspace{12pt}
\end{center}
}]

% ---- Abstract -------------------------------------------------------
\noindent\textbf{\textit{Abstract}}
\newline
\textit{%
This report presents an experimental study of convolutional neural networks
for image classification on a CIFAR-10 subset with 5 classes. I implemented a
SmallCNN from scratch (70.30\% test accuracy), compared ResNet-18 feature
extraction (86.40\%) versus fine-tuning (93.70\%), and evaluated data
augmentation strategies including Mixup. The best performance was achieved by
fine-tuning ResNet-18 with no augmentation (94.00\%), demonstrating that
pretrained features can be highly effective on small datasets. The key lesson
is that transfer learning with fine-tuning dramatically outperforms training
from scratch on limited data, and augmentation strategies should be chosen
carefully as they can sometimes hurt performance.}

\vspace{4pt}
\noindent\textbf{\textit{Index Terms}}---convolutional neural networks,
transfer learning, data augmentation, image classification, PyTorch.

\vspace{6pt}

% =====================================================================
\section{Introduction}
Convolutional neural networks are well-suited for vision tasks due to two key
properties: parameter sharing (the same filter is applied across all spatial
locations) and local connectivity (each neuron connects only to a local region
of the input). This report covers a complete experimental study of CNNs on a
CIFAR-10 subset, including a small CNN trained from scratch, transfer learning
with ResNet-18, and a data-augmentation study.

This work maps directly to the homework: Section~\ref{sec:theory} presents a
theoretical analysis (Part A), Sections~\ref{sec:setup} and~\ref{sec:results}
describe and evaluate the experiments (Part B), and Section~\ref{sec:discussion}
explains the implementation details (Part C).

\noindent\textit{Contributions:}
\begin{itemize}[leftmargin=1.2em,itemsep=1pt,topsep=2pt]
  \item a \texttt{SmallCNN} baseline trained on a CIFAR-10 subset;
  \item a comparison of ResNet-18 \emph{feature extraction} vs.\ \emph{fine-tuning};
  \item an augmentation study including self-implemented Mixup.
\end{itemize}

% =====================================================================
\section{Background and Theoretical Analysis}\label{sec:theory}

Throughout we use the convolution output-size formula
\begin{equation}\label{eq:convsize}
  n_{\text{out}} = \left\lfloor \frac{n + 2P - k}{s} \right\rfloor + 1 .
\end{equation}

\subsection{Convolution output size and parameter count}

An RGB input of shape $3\times64\times64$ passes through a convolution with
$C_{\text{out}}=16$, $k=5$, $s=1$, $P=2$, then a $2\times2$ max-pool ($s=2$).

\textbf{(a)} Convolution:
\begin{equation}
n_{\text{out}} = \left\lfloor \frac{64 + 2(2) - 5}{1} \right\rfloor + 1
= \lfloor 63 \rfloor + 1 = 64
\end{equation}
Shape after conv: $16 \times 64 \times 64$.

Max-pool: $n_{\text{out}} = \lfloor (64 - 2)/2 \rfloor + 1 = 31 + 1 = 32$.
Shape after pool: $16 \times 32 \times 32$.

\textbf{(b)} Learnable parameters:
\begin{equation}
(k^2 \cdot C_{\text{in}} + 1) \cdot C_{\text{out}}
= (5^2 \cdot 3 + 1) \cdot 16
= (75 + 1) \cdot 16
= 1,216
\end{equation}

\textbf{(c)} Dense layer: flatten $3\times64\times64 = 12,288$ to $16$ outputs:
$12,288 \times 16 + 16 = 196,624$ parameters.
Ratio: $196,624 / 1,216 \approx 161.7$.

The two convolution properties are: (1) \textbf{Parameter sharing}: same
filter weights applied across all spatial positions; (2) \textbf{Local
connectivity}: each output neuron connects to only $k\times k\times C_{\text{in}}$
inputs, not all.

\subsection{Receptive field and dilation}

For three stacked $3\times3$ convolutions, stride 1, no dilation:

\textbf{(a)} Using $r_\ell = r_{\ell-1} + (k_\ell - 1)\prod_{i<\ell} s_i$:
\begin{align*}
r_1 &= 1 + (3-1)\cdot 1 = 3 \\
r_2 &= 3 + (3-1)\cdot 1 = 5 \\
r_3 &= 5 + (3-1)\cdot 1 = 7
\end{align*}
Receptive field after third layer: $7 \times 7$.

\textbf{(b)} With second layer dilation $d=2$:
Effective kernel size: $k_{\text{eff}} = k + (k-1)(d-1) = 3 + 2(1) = 5$.
\begin{align*}
r_1 &= 3 \\
r_2 &= 3 + (5-1)\cdot 1 = 7 \\
r_3 &= 7 + (3-1)\cdot 1 = 9
\end{align*}
Receptive field: $9 \times 9$.

\textbf{(c)} Parameter count (per channel-pair):
- Three $3\times3$ layers: $3 \times 9 = 27$
- One $7\times7$ layer: $49$

Two reasons to prefer the stack: (1) fewer parameters ($27$ vs $49$), and
(2) more non-linearity (three ReLU activations vs one).

\subsection{Transposed convolution for a decoder}

A decoder upsamples $8\times8 \to 16\times16$ using:
\begin{equation}
n_{\text{out}} = s(n_{\text{in}} - 1) + k - 2P
\end{equation}

\textbf{(a)} Choose $s=2$: $16 = 2(8-1) + k - 2P = 14 + k - 2P$.
Thus $k - 2P = 2$. Choose $k=4, P=1$: $4 - 2(1) = 2$. Valid $(k,s,P) = (4,2,1)$.

\textbf{(b)} Condition to avoid checkerboard artifacts: $k$ must be divisible by
$s$ ($k \% s = 0$). My choice $k=4, s=2$ satisfies this. An alternative to
transposed convolution is nearest-neighbor upsampling followed by a regular
convolution.

\subsection{Detection metrics: IoU, precision, recall}

For boxes $A=(1,1,5,4)$ and $B=(3,2,7,6)$:

\textbf{(a)} Intersection: $x_1=\max(1,3)=3$, $y_1=\max(1,2)=2$,
$x_2=\min(5,7)=5$, $y_2=\min(4,6)=4$.
Area: $(5-3)(4-2) = 2 \times 2 = 4$.

Union: $12 + 16 - 4 = 24$.

$\mathrm{IoU} = 4/24 = 1/6 \approx 0.1667$.

\textbf{(b)} Precision/Recall Table:

\begin{table}[t]
\caption{Cumulative precision/recall}
\label{tab:pr}
\centering
\begin{tabular}{cccc}
\toprule
\# & TP/FP & Cum.\ precision & Cum.\ recall \\
\midrule
1 & TP & 1/1 = 1.000 & 1/4 = 0.250 \\
2 & FP & 1/2 = 0.500 & 1/4 = 0.250 \\
3 & TP & 2/3 = 0.667 & 2/4 = 0.500 \\
4 & TP & 3/4 = 0.750 & 3/4 = 0.750 \\
5 & FP & 3/5 = 0.600 & 3/4 = 0.750 \\
\bottomrule
\end{tabular}
\end{table}

\textbf{(c)} Average Precision is the area under the precision-recall curve,
computed via interpolation at each recall level. AP is preferred over a single
precision/recall pair because it summarizes performance across all confidence
thresholds, rather than depending on an arbitrary threshold choice.

\subsection{Residual skip connections}

\textbf{(a)} The degradation problem: As network depth increases, training
accuracy saturates and then degrades rapidly — this is NOT overfitting, as
deeper networks have higher training error, indicating the optimization problem
is harder.

\textbf{(b)} For $\mathbf{y} = F(\mathbf{x}) + \mathbf{x}$:
\begin{equation}
\frac{\partial \mathbf{y}}{\partial \mathbf{x}}
= \frac{\partial F}{\partial \mathbf{x}} + \mathbf{I}
\end{equation}
Even when $\partial F/\partial \mathbf{x}$ is tiny (near zero), the $+\mathbf{I}$
term ensures gradients flow through the skip connection, allowing early layers
to receive signal in very deep networks.

\textbf{(c)} Driving $F \to 0$ (learning no change) is easier than learning an
identity map through plain layers because plain layers have no built-in
mechanism to bypass transformations.

\subsection{Choosing a transfer-learning strategy}

\textbf{(i)} 300 labelled images, similar to ImageNet: \textbf{Feature
extraction} — freeze the backbone, train only a new head (limited data, but
features are relevant).

\textbf{(ii)} 200k labelled images, similar classes: \textbf{Fine-tuning}
— with abundant data, fine-tune the whole network to adapt domain-specific
patterns without overfitting.

\textbf{(iii)} 500 grayscale medical scans, very different: \textbf{Partial
fine-tuning} — freeze early layers (general features), fine-tune later layers
or use feature extraction with a deeper head.

\textbf{(a)} The backbone's pretrained weights were optimized for specific input
statistics (mean, std, and spatial size). Changing the input distribution
changes what the weights see, effectively breaking the pretraining.

\textbf{(b)} Pretrained weights are already good. A large learning rate would
make large updates that destroy these useful features. A smaller LR allows
gentle adaptation without catastrophic forgetting.

% =====================================================================
\section{Experimental Setup}\label{sec:setup}

\subsection{Data}
I used a CIFAR-10 subset with 5 classes (airplane, automobile, bird, cat, deer),
with 800 train and 200 test images per class, seeded by Student ID 121314. A
validation set of 160 images per class was also extracted from the training
data. For ResNet, images were resized to $224\times224$ and normalized with
ImageNet mean/std (mean $[0.485,0.456,0.406]$, std $[0.229,0.224,0.225]$) as
required by the pretrained backbone (Problem 6a).

\subsection{Models}

\texttt{SmallCNN} consists of 3 conv blocks (Conv--BatchNorm--ReLU--MaxPool),
global average pooling, and a linear head. The architecture is summarized in
Table~\ref{tab:arch}. For transfer learning, \texttt{build\_resnet18} provides
two modes: feature extraction (backbone frozen, new head only, 2,565 trainable
params) and fine-tuning (whole network trainable, 11.2M params, smaller LR).

\begin{table}[t]
\caption{SmallCNN architecture}
\label{tab:arch}
\centering
\begin{tabular}{lll}
\toprule
Stage & Layer & Output shape \\
\midrule
Input   & --                                   & $3\times32\times32$ \\
Block 1 & Conv(3$\to$32)--BN--ReLU--Pool(2) & $32\times16\times16$ \\
Block 2 & Conv(32$\to$64)--BN--ReLU--Pool(2)& $64\times8\times8$ \\
Block 3 & Conv(64$\to$128)--BN--ReLU--Pool(2)& $128\times4\times4$ \\
Head    & GAP $+$ Linear(128$\to$5)           & $5$ \\
\bottomrule
\end{tabular}
\end{table}

\subsection{Augmentation}

Mixup forms convex combinations:
\begin{align}
  \tilde{\mathbf{x}} &= \lambda\,\mathbf{x}_i + (1-\lambda)\,\mathbf{x}_j, \\
  \mathcal{L} &= \lambda\,\ell(\hat{\mathbf{y}}, y_i)
               + (1-\lambda)\,\ell(\hat{\mathbf{y}}, y_j),
\end{align}
with $\lambda \sim \mathrm{Beta}(1,1)$. I implemented Mixup (CutMix was also
implemented but Mixup was used for the augmentation study as specified in B3).

\subsection{Training and reproducibility}

All experiments used Adam optimizer with cross-entropy loss, batch size 64, and
20 epochs. Learning rates were $10^{-3}$ for SmallCNN and feature extraction,
and $10^{-4}$ for fine-tuning (Problem 6b). All randomness was fixed with
Student ID 121314 as the seed, controlling data split, weight initialization,
and shuffling, ensuring every figure regenerates via \texttt{run\_all.py}.

% =====================================================================
\section{Results}\label{sec:results}

\subsection{SmallCNN baseline (B1)}

SmallCNN achieved a test accuracy of 70.30\%. As shown in Fig.~\ref{fig:cnn},
the model converges steadily, but the gap between training and validation
accuracy (88.06\% vs 73.38\%) indicates overfitting — the model memorizes
training data but fails to generalize perfectly to unseen data.

\begin{figure}[t]
\centering
\includegraphics[width=\linewidth]{cnn_curves.png}
\caption{Train/validation loss and accuracy vs.\ epoch for SmallCNN.}
\label{fig:cnn}
\end{figure}

\subsection{Transfer learning: feature extraction vs.\ fine-tuning (B2)}

Table~\ref{tab:transfer} shows that fine-tuning significantly outperforms
feature extraction (93.70\% vs 86.40\%) on this small dataset. The trainable
parameter counts differ dramatically (2,565 vs 11.2M), but the additional
capacity of fine-tuning is beneficial. As shown in Fig.~\ref{fig:transfer},
fine-tuning converges faster and achieves higher validation accuracy. This
matches the transfer-learning guidance: with a relatively small but
representative dataset, fine-tuning adapts the pretrained features to the
specific domain more effectively than freezing them.

\begin{table}[t]
\caption{Transfer-learning comparison}
\label{tab:transfer}
\centering
\begin{tabular}{lrr}
\toprule
Mode & Trainable params & Test acc. \\
\midrule
Feature extraction (frozen) & 2,565      & 0.8640 \\
Fine-tuning (whole net)     & 11,179,077 & 0.9370 \\
\bottomrule
\end{tabular}
\end{table}

\begin{figure}[t]
\centering
\includegraphics[width=\linewidth]{transfer_compare.png}
\caption{Validation accuracy: feature extraction, fine-tuning, and the
SmallCNN baseline, on one axis.}
\label{fig:transfer}
\end{figure}

\subsection{Data-augmentation study (B3)}

Using the best architecture (fine-tuned ResNet-18), I trained under three
regimes. Table~\ref{tab:aug} and Fig.~\ref{fig:aug} show the results.

\begin{table}[t]
\caption{Augmentation study}
\label{tab:aug}
\centering
\begin{tabular}{lr}
\toprule
Regime & Test acc. \\
\midrule
(a) None                        & 0.9400 \\
(b) Standard (crop $+$ flip)    & 0.9170 \\
(c) Mixup                       & 0.9340 \\
\bottomrule
\end{tabular}
\end{table}

Surprisingly, no augmentation achieved the highest test accuracy (94.00\%).
Standard augmentation (RandomCrop + HorizontalFlip) performed worst (91.70\%),
while Mixup (93.40\%) fell between the two. This suggests that with only 3,200
training images and a strong pretrained ResNet-18 backbone, augmentation may
introduce unnecessary distribution shift rather than improving generalization.
The validation curves confirm that "No Augmentation" consistently outperforms
or matches the other regimes across all epochs, indicating a stable effect.

\begin{figure}[t]
\centering
\includegraphics[width=\linewidth]{augment_compare.png}
\caption{Validation accuracy across the three augmentation regimes.}
\label{fig:aug}
\end{figure}

\subsection{Summary and discussion (B4)}

\begin{table}[t]
\caption{Overall results}
\label{tab:summary}
\centering
\begin{tabular}{llr}
\toprule
Model & Setting & Test acc. \\
\midrule
SmallCNN  & from scratch       & 0.7030 \\
ResNet-18 & feature extraction & 0.8640 \\
ResNet-18 & fine-tuning        & 0.9370 \\
ResNet-18 & No Augmentation    & 0.9400 \\
\bottomrule
\end{tabular}
\end{table}

Transfer learning dramatically outperforms training from scratch (93.70\% vs
70.30\%) on this small dataset. Among transfer strategies, fine-tuning
outperforms feature extraction. The augmentation study reveals that on this
dataset, the pretrained features are already robust enough that no augmentation
performed best — standard augmentation actually hurt performance, while Mixup
provided a middle ground. The best overall setting was fine-tuned ResNet-18
with no augmentation (94.00\%).

% =====================================================================
\section{Implementation Discussion}\label{sec:discussion}

\subsection{SmallCNN and global average pooling (C1)}

My SmallCNN layers are:
\begin{itemize}
\item Block 1: Conv2d(3, 32, k=3, p=1) $\to$ BatchNorm $\to$ ReLU $\to$ MaxPool(2)
\item Block 2: Conv2d(32, 64, k=3, p=1) $\to$ BatchNorm $\to$ ReLU $\to$ MaxPool(2)
\item Block 3: Conv2d(64, 128, k=3, p=1) $\to$ BatchNorm $\to$ ReLU $\to$ MaxPool(2)
\item Head: AdaptiveAvgPool2d(1) $\to$ Flatten $\to$ Linear(128, 5)
\end{itemize}

Using \eqref{eq:convsize}:
- Block 1: $32 \to (32-2)/2 + 1 = 16$
- Block 2: $16 \to (16-2)/2 + 1 = 8$
- Block 3: $8 \to (8-2)/2 + 1 = 4$

Spatial size at GAP input: $4 \times 4$.

I used GAP instead of flatten+dense because it reduces parameters from
$128 \times 4 \times 4 \times 5 = 10,240$ to $128 \times 5 = 640$ (93.7\%
reduction), significantly reducing overfitting on the small dataset.

\subsection{Freezing, concretely (C2)}

In feature-extraction mode, I set:
\begin{lstlisting}
for param in model.parameters():
    param.requires_grad = False
\end{lstlisting}
This freezes all ResNet-18 backbone parameters. The new `model.fc` layer is
created after this loop, so it has `requires_grad=True` by default.

Parameter counts (from `count_trainable_params()`):
- Feature extraction: 2,565 (only the new Linear(512, 5) head)
- Fine-tuning: 11,179,077 (entire network)

\subsection{Mixup/CutMix in code (C3)}

Mixup implementation from `augment.py`:
\begin{lstlisting}
lam = np.random.beta(alpha, alpha)
idx = torch.randperm(batch_size)
x_mixed = lam * x + (1 - lam) * x[idx]
\end{lstlisting}

Loss combination:
\begin{lstlisting}
def mix_criterion(loss_fn, logits, y_a, y_b, lam):
    return lam * loss_fn(logits, y_a) + (1 - lam) * loss_fn(logits, y_b)
\end{lstlisting}

The loss is a convex combination because $\lambda + (1-\lambda) = 1$. A single
cross-entropy would only work for one label; the convex combination properly
handles the mixed supervision since the input is a mix of two images.

\subsection{A pitfall you hit (C4)}

The most significant bug was in B3's test evaluation. Originally, I used
`load_base_model()` to load the B2 checkpoint for all three regimes,
resulting in identical test accuracies (0.9370, 0.9370, 0.9370). The fix was
to save each regime's trained model to `augment_{regime}.pth` and load them
with `load_regime_model(regime, device)` instead of the shared B2 base model.
After the fix, test accuracies differed (0.9400, 0.9170, 0.9340).

\subsection{Reproducibility (C5)}

My Student ID is 121314. The seed flows through:
1. **Data split:** `np.random.RandomState(seed)` controls shuffling and
   stratified sampling.
2. **Weight initialization:** PyTorch's default initialization is
   deterministic when `torch.manual_seed(seed)` is set.
3. **DataLoader shuffling:** Uses PyTorch's RNG, seeded via
   `torch.manual_seed(seed)`.
4. **CUDA:** `torch.cuda.manual_seed_all(seed)` ensures GPU determinism.

Running `python run_all.py` twice produces identical figures.

% =====================================================================
\section{Pretrained-Model Inference (Bonus)}

I ran pretrained torchvision models on 3 photos of soccer players (Ngolo Kante,
Thomas Muller, and a match scene). The annotated outputs are saved as
`figures/zoo_*.png`.

\textbf{Detection (Faster R-CNN with FPN):}
FPN (Feature Pyramid Network) uses a top-down pathway with lateral connections
to build high-resolution semantic features at all scales, enabling detection of
both small and large objects. NMS removes duplicate detections by keeping only
the highest-confidence box per object.

\textbf{Segmentation (Mask R-CNN with RoIAlign):}
RoIPool uses quantization (floor/ceil) causing misalignment between the RoI and
feature map. RoIAlign uses bilinear interpolation to maintain sub-pixel
precision, critical for accurate mask prediction.

\textbf{Pose (Keypoint R-CNN with heatmap regression):}
Direct (x,y) regression is unstable because the loss landscape is non-convex
with many local minima. Heatmap regression produces a 2D probability
distribution that provides spatial context and is more robust to occlusions.

% =====================================================================
\section{Conclusion}

The best performance on this CIFAR-10 subset was achieved by fine-tuning
ResNet-18 without augmentation (94.00\% test accuracy). Transfer learning
proved dramatically superior to training from scratch (93.70\% vs 70.30\%).
The augmentation study showed that on this small dataset with a strong
pretrained backbone, standard augmentation can hurt performance, while Mixup
provides a middle ground. The key takeaway is that transfer learning with
carefully chosen hyperparameters is essential for small datasets, and
augmentation strategies should be validated rather than assumed beneficial.

% =====================================================================
\section*{Acknowledgment: use of AI tools}

Used Claude (Anthropic) for:
- Code structure and architecture planning
- Debugging the B3 model checkpoint bug
- Explaining transposed convolution sizing
- Generating the report template

All code was reviewed and understood before submission.

% =====================================================================
\begin{thebibliography}{10}
\bibitem{lecun} Y.~LeCun, L.~Bottou, Y.~Bengio, and P.~Haffner, ``Gradient-based learning applied to document recognition,'' \emph{Proc.\ IEEE}, vol.~86, no.~11, pp.~2278--2324, 1998.
\bibitem{alexnet} A.~Krizhevsky, I.~Sutskever, and G.~E.~Hinton, ``ImageNet classification with deep convolutional neural networks,'' in \emph{NeurIPS}, 2012.
\bibitem{vgg} K.~Simonyan and A.~Zisserman, ``Very deep convolutional networks for large-scale image recognition,'' in \emph{ICLR}, 2015.
\bibitem{resnet} K.~He, X.~Zhang, S.~Ren, and J.~Sun, ``Deep residual learning for image recognition,'' in \emph{CVPR}, 2016, pp.~770--778.
\bibitem{bn} S.~Ioffe and C.~Szegedy, ``Batch normalization: Accelerating deep network training by reducing internal covariate shift,'' in \emph{ICML}, 2015.
\bibitem{mixup} H.~Zhang, M.~Cisse, Y.~N.~Dauphin, and D.~Lopez-Paz, ``mixup: Beyond empirical risk minimization,'' in \emph{ICLR}, 2018.
\bibitem{cutmix} S.~Yun \emph{et al.}, ``CutMix: Regularization strategy to train strong classifiers with localizable features,'' in \emph{ICCV}, 2019.
\bibitem{fasterrcnn} S.~Ren, K.~He, R.~Girshick, and J.~Sun, ``Faster R-CNN: Towards real-time object detection with region proposal networks,'' in \emph{NeurIPS}, 2015.
\bibitem{maskrcnn} K.~He, G.~Gkioxari, P.~Doll\'ar, and R.~Girshick, ``Mask R-CNN,'' in \emph{ICCV}, 2017.
\bibitem{pytorch} A.~Paszke \emph{et al.}, ``PyTorch: An imperative style, high-performance deep learning library,'' in \emph{NeurIPS}, 2019.
\end{thebibliography}

\end{document}